\documentclass[11pt,a4paper]{article}
\usepackage[utf8]{inputenc}
\usepackage[T1]{fontenc}
\usepackage{lmodern}
\usepackage[portuguese]{babel}
\usepackage[a4paper,margin=2.5cm]{geometry}
\usepackage{xcolor}
\definecolor{TextEspresso}{HTML}{4A4238}
\definecolor{NudeTaupe}{HTML}{BFAFA6}
\definecolor{SoftBlush}{HTML}{D9C5B2}
\usepackage{titlesec}
\usepackage{enumitem}
\usepackage{listings}
\usepackage{hyperref}
\hypersetup{colorlinks=true, linkcolor=TextEspresso, urlcolor=TextEspresso}

\titleformat{\section}{\color{TextEspresso}\Large\bfseries}{\thesection}{0.5em}{}[{\color{NudeTaupe}\titlerule[0.8pt]}]
\titleformat{\subsection}{\color{TextEspresso}\large\bfseries}{\thesubsection}{0.5em}{}

\lstset{
  basicstyle=\ttfamily\small,
  breaklines=true,
  frame=single,
  rulecolor=\color{NudeTaupe},
  columns=fullflexible,
  showstringspaces=false
}

\newcommand{\guiaotitle}[2]{%
  \begin{center}
  {\Large\bfseries\color{TextEspresso} #1}\\[0.3cm]
  {\normalsize\color{NudeTaupe} #2}\\[0.2cm]
  {\color{NudeTaupe}\rule{4cm}{0.5pt}}
  \end{center}
  \vspace{0.5cm}
}

\begin{document}
\guiaotitle{Guião 06 -- SDN com Mininet e Open vSwitch}{Infraestruturas de Computação na Nuvem, Aula 06}

\section{Objetivos}
\begin{itemize}[itemsep=0.2cm]
  \item Observar a rede virtual real do servidor Proxmox (bridges) e do cluster Kubernetes (overlay CNI).
  \item Criar uma topologia de rede virtual com Mininet, usando switches Open vSwitch (OVS).
  \item Observar o comportamento da rede sem controller e com um controller OpenFlow ligado.
  \item Inspecionar a tabela de fluxos instalada num switch OVS, incluindo prioridades, contadores e timeouts.
  \item Implementar uma política de rede simples (bloqueio de tráfego entre dois hosts) através do controller.
\end{itemize}

\section{Pré-requisitos e Material}
\begin{itemize}[itemsep=0.2cm]
  \item A VM \texttt{icn-lab} no Proxmox (Guião 01), onde será instalado o Mininet. Recomenda-se pelo menos 2~GB de RAM.
  \item Acesso à interface web do Proxmox, para a observação inicial da rede do servidor.
  \item Acesso ao cluster Kubernetes da Aula 05, se ainda estiver montado (opcional, para a Parte 1).
  \item Acesso root/sudo na VM: o Mininet cria interfaces de rede virtuais e exige privilégios.
\end{itemize}

\textbf{Nota:} o Mininet corre \emph{dentro} da sua VM, criando aí uma rede virtual emulada. Não interfere com a rede do servidor Proxmox nem com a das restantes VMs da turma.

\section{Introdução}
Esta aula tem duas partes complementares. Na primeira, observamos redes virtuais \emph{reais} já em funcionamento: as bridges que o Proxmox usa para ligar as VMs da turma, e a rede overlay que o plugin CNI criou no cluster Kubernetes da Aula 05. Na segunda, construímos do zero uma rede definida por software com Mininet e um controller OpenFlow, para ver o mecanismo por dentro.

O Mininet emula hosts, switches e ligações num único Linux, usando os mesmos namespaces de rede estudados na Aula 04. Os switches são instâncias de Open vSwitch -- o mesmo software usado em datacenters reais. Vamos primeiro observar que, sem controller, o switch não sabe o que fazer aos pacotes; depois ligamos um controller Ryu que instala as regras na tabela de fluxos; e por fim modificamo-lo para aplicar uma política de rede.

\section{Parte 1 -- Observar redes virtuais reais}
\begin{enumerate}[label=\arabic*.]

  \item \textbf{Observar as bridges do Proxmox.} Na interface web, selecione o \emph{node} e abra \emph{System} $\rightarrow$ \emph{Network}. Identifique a bridge (tipicamente \texttt{vmbr0}) à qual as VMs estão ligadas e registe a que interface física corresponde.

  Se tiver acesso à shell do node, observe a mesma informação por CLI:
\begin{lstlisting}
$ ip link show type bridge
$ brctl show 2>/dev/null || bridge link
\end{lstlisting}
  Esta bridge é um switch virtual em software: desempenha, para as VMs, o papel que um switch físico desempenha para servidores físicos.

  \item \textbf{Observar a rede overlay do Kubernetes} (se o cluster da Aula 05 ainda estiver disponível). A partir do control plane:
\begin{lstlisting}
$ kubectl get pods -o wide -A | head -20
$ ip addr show
\end{lstlisting}
  Procure interfaces criadas pelo plugin CNI (ex.\ \texttt{flannel.1}, \texttt{cni0}). A interface \texttt{flannel.1} é um \textbf{VTEP}: o ponto onde o tráfego entre Pods de nós diferentes é encapsulado em VXLAN.

  Compare a sub-rede dos Pods com a sub-rede das VMs:
\begin{lstlisting}
$ kubectl get nodes -o jsonpath='{.items[*].spec.podCIDR}'
\end{lstlisting}
  Registe as duas gamas de endereços. A dos Pods é a rede \emph{overlay}; a das VMs é o \emph{underlay}.

\end{enumerate}

\section{Parte 2 -- Construir uma rede SDN com Mininet}
\begin{enumerate}[label=\arabic*.,resume]

  \item \textbf{Instalar o Mininet e o Ryu} na VM \texttt{icn-lab}:
\begin{lstlisting}
$ sudo apt update
$ sudo apt install -y mininet openvswitch-switch python3-pip
$ sudo systemctl enable --now openvswitch-switch
\end{lstlisting}
  O controller Ryu tem dependências sensíveis à versão do Python; instale-o num ambiente virtual para não interferir com o sistema:
\begin{lstlisting}
$ sudo apt install -y python3-venv
$ python3 -m venv ~/ryu-env
$ ~/ryu-env/bin/pip install ryu
\end{lstlisting}
  \emph{Nota:} caso a instalação do Ryu falhe por incompatibilidade de versões (é um projeto pouco mantido), a docente indicará a alternativa a usar -- tipicamente uma versão fixada do pacote, ou o controller de referência do próprio Mininet para os passos que não exigem código próprio.

  \item Confirmar que o Mininet e o Open vSwitch estão operacionais:
\begin{lstlisting}
$ sudo mn --version
$ sudo ovs-vsctl --version
\end{lstlisting}

  \item Criar uma topologia simples com três hosts ligados a um único switch, sem indicar controller (o Mininet arranca um controller de referência por omissão):
\begin{lstlisting}
$ sudo mn --topo single,3 --mac --switch ovsk
\end{lstlisting}
  Dentro da consola do Mininet, testar a conetividade:
\begin{lstlisting}
mininet> pingall
\end{lstlisting}
  Deve funcionar, pois o controller de referência do Mininet instala automaticamente um comportamento de \emph{learning switch}.

  \item Sair do Mininet (\texttt{exit}) e repetir a topologia, desta vez sem qualquer controller, para observar o efeito de um switch órfão:
\begin{lstlisting}
$ sudo mn --topo single,3 --mac --switch ovsk --controller none
mininet> pingall
\end{lstlisting}
  Sem controller, o switch não tem forma de aprender para onde encaminhar os pacotes: os pings falham. Isto ilustra a dependência do plano de dados em relação ao plano de controlo.

  \item Sair do Mininet. Num terminal separado, criar um controller simples baseado em Ryu, ficheiro \texttt{learning\_switch.py}:
\begin{lstlisting}
from ryu.base import app_manager
from ryu.controller import ofp_event
from ryu.controller.handler import MAIN_DISPATCHER, CONFIG_DISPATCHER, set_ev_cls
from ryu.ofproto import ofproto_v1_3
from ryu.lib.packet import packet, ethernet

class LearningSwitch(app_manager.RyuApp):
    OFP_VERSIONS = [ofproto_v1_3.OFP_VERSION]

    def __init__(self, *args, **kwargs):
        super(LearningSwitch, self).__init__(*args, **kwargs)
        self.mac_to_port = {}

    @set_ev_cls(ofp_event.EventOFPSwitchFeatures, CONFIG_DISPATCHER)
    def switch_features_handler(self, ev):
        datapath = ev.msg.datapath
        parser = datapath.ofproto_parser
        ofproto = datapath.ofproto
        match = parser.OFPMatch()
        actions = [parser.OFPActionOutput(ofproto.OFPP_CONTROLLER,
                                           ofproto.OFPCML_NO_BUFFER)]
        self.add_flow(datapath, 0, match, actions)

    def add_flow(self, datapath, priority, match, actions):
        parser = datapath.ofproto_parser
        ofproto = datapath.ofproto
        inst = [parser.OFPInstructionActions(ofproto.OFPIT_APPLY_ACTIONS, actions)]
        mod = parser.OFPFlowMod(datapath=datapath, priority=priority,
                                 match=match, instructions=inst)
        datapath.send_msg(mod)

    @set_ev_cls(ofp_event.EventOFPPacketIn, MAIN_DISPATCHER)
    def packet_in_handler(self, ev):
        msg = ev.msg
        datapath = msg.datapath
        ofproto = datapath.ofproto
        parser = datapath.ofproto_parser
        in_port = msg.match['in_port']

        pkt = packet.Packet(msg.data)
        eth = pkt.get_protocols(ethernet.ethernet)[0]
        dst = eth.dst
        src = eth.src
        dpid = datapath.id
        self.mac_to_port.setdefault(dpid, {})
        self.mac_to_port[dpid][src] = in_port

        out_port = self.mac_to_port[dpid].get(dst, ofproto.OFPP_FLOOD)
        actions = [parser.OFPActionOutput(out_port)]

        if out_port != ofproto.OFPP_FLOOD:
            match = parser.OFPMatch(in_port=in_port, eth_dst=dst)
            self.add_flow(datapath, 1, match, actions)

        data = msg.data if msg.buffer_id == ofproto.OFP_NO_BUFFER else None
        out = parser.OFPPacketOut(datapath=datapath, buffer_id=msg.buffer_id,
                                   in_port=in_port, actions=actions, data=data)
        datapath.send_msg(out)
\end{lstlisting}

  \item Arrancar o controller num terminal (usando o ambiente virtual criado):
\begin{lstlisting}
$ ~/ryu-env/bin/ryu-manager learning_switch.py
\end{lstlisting}

  \item Noutro terminal, arrancar o Mininet a apontar para este controller externo (por omissão, escuta em \texttt{127.0.0.1:6633}):
\begin{lstlisting}
$ sudo mn --topo single,3 --mac --switch ovsk \
    --controller remote,ip=127.0.0.1,port=6633
mininet> pingall
\end{lstlisting}
  Confirmar que a conetividade funciona novamente, agora com as regras a serem instaladas pelo nosso controller.

  \item Inspecionar a tabela de fluxos instalada no switch (executar num terminal fora do Mininet, com a topologia ainda a correr):
\begin{lstlisting}
$ sudo ovs-ofctl dump-flows s1
\end{lstlisting}
  Identificar as entradas \emph{match -- action} correspondentes aos pares de hosts que já comunicaram. Para cada entrada, localize também os elementos discutidos na teoria:
\begin{lstlisting}
- priority=...    (qual regra vence quando varias correspondem)
- n_packets=...   (contador de pacotes que corresponderam)
- n_bytes=...     (contador de bytes)
- duration=...    (ha quanto tempo a entrada existe)
\end{lstlisting}
  Gere mais tráfego e volte a consultar a tabela, confirmando que os contadores aumentam:
\begin{lstlisting}
mininet> h1 ping -c 5 h2
$ sudo ovs-ofctl dump-flows s1
\end{lstlisting}

  \item Modificar o controller para bloquear especificamente o tráfego entre \texttt{h1} e \texttt{h3}, mantendo a comunicação com \texttt{h2}: no método \texttt{packet\_in\_handler}, antes de calcular \texttt{out\_port}, adicionar uma verificação sobre os endereços MAC de origem/destino conhecidos de \texttt{h1} e \texttt{h3} (obter os MACs com \texttt{h1 ifconfig} e \texttt{h3 ifconfig} dentro da consola do Mininet) e, nesse caso, instalar uma regra com lista de \texttt{actions} vazia (drop) em vez de reencaminhar.

  \item Reiniciar o controller e o Mininet com a alteração, e confirmar o novo comportamento:
\begin{lstlisting}
mininet> h1 ping -c 3 h2
mininet> h1 ping -c 3 h3
\end{lstlisting}

\end{enumerate}

\section{Entregável / Questões de Verificação}
\begin{enumerate}[label=\alph*)]
  \setlength\itemsep{0.2cm}
  \item Enviar o ficheiro \texttt{learning\_switch.py} final (com o bloqueio h1--h3) e a saída de \texttt{ovs-ofctl dump-flows s1} antes e depois da alteração.
  \item Que bridge do Proxmox liga as VMs da turma, e a que interface física corresponde? Em que sentido é ela um "switch virtual"?
  \item Registe a sub-rede dos Pods e a sub-rede das VMs observadas na Parte 1. Qual é o \emph{overlay} e qual é o \emph{underlay}? Que interface faz o papel de VTEP?
  \item Porque é que o \texttt{pingall} falha quando não existe nenhum controller ligado ao switch? O que isto demonstra sobre a relação entre plano de controlo e plano de dados?
  \item Que informação usa o controller para decidir a porta de saída de um pacote (\texttt{mac\_to\_port})? Que camada do modelo OSI está a ser usada nessa decisão?
  \item O controller implementado neste guião segue um modelo reativo ou proativo? Justifique com base no que acontece quando chega um pacote sem regra correspondente.
  \item Observou os contadores (\texttt{n\_packets}) a aumentar nas entradas de fluxo. Para que serve, na prática, esta informação num datacenter real?
  \item Em que consiste, neste guião, o plano de dados e em que consiste o plano de controlo? Identifique claramente os dois componentes usados.
\end{enumerate}

\end{document}
